\begin{afterword}
\summaryhead

The author says that comments on the previous posts have clarified Moore's paradox. The commenter jimrandomh suggested that many people cannot tell ``I believe X'' from ``X.'' The author disagrees, since young children understand false belief, but proposes a similar idea: many people may not distinguish believing something from endorsing it. The word ``belief'' has both senses, as in ``I believe in democracy.''

So the woman of the earlier posts found reasons to think it good to believe that people are nice, felt warmly toward that belief, and took the warmth for belief itself. Seeing that people are not so nice, she said ``I believe people are nicer than they are.'' No one is taught what real belief feels like: a believed statement simply seems like the way the world is.

The commenter Kurige, who believes both that God instilled a sense of right and wrong and that morality evolved, is diagnosed the same way. Kurige is asked to state both claims without ``I believe'' and to notice that they are harder to swallow.

\responsehead

A post titled ``Moore's Paradox'' never says what the paradox is. It is the puzzle of why a sentence like ``It is raining, but I do not believe that it is raining'' seems absurd to assert, though it is consistent and may be true. A reader who does not already know this meets the name without the problem, and an example appears only in a parenthesis in the next post.

The explanation the post offers does not address the paradox either. If the woman used ``believe'' to mean ``endorse,'' her sentence is not a Moorean sentence at all: ``I endorse believing that people are nice, but they are not'' is neither inconsistent nor odd. The post explains how someone can utter a sentence of that form while meaning something else. That may be a real mechanism, and its distinction between believing and endorsing is sound; Merriam-Webster records both senses of ``believe.'' But it does not answer the question philosophers ask, which is why the sentence is absurd when ``believe'' means believe.

The mechanism is also stated as fact about a mind the author knows from one conversation: ``she mistook the positive affect attached to the quoted belief.'' Three days earlier the same sentence of hers was explained by her belief that she had deceived herself. The post calls the new account ``another mechanism'' and does not say how to tell which one applies.

The test proposed to Kurige does not separate the two states it is meant to separate. Dropping ``I believe'' also drops a hedge, so the flat sentence may be harder to say for anyone, whatever they believe. Kurige replied on this post with an explanation of this kind, comparing the difference to a scientist who will say ``according to our model, the higgs-boson should exist'' more readily than ``the higgs-boson exists.'' The diagnosis of Kurige also supplies reasons, good communities to join and good feelings, that Kurige's comment did not give.

In short: a post named after a paradox it never states, whose explanation, if right, shows that its case is not an instance of the paradox, and whose proposed self-test does not separate a hedge from a lack of belief.
\end{afterword}
